\documentclass[12pt]{amsart}
\usepackage{amsmath,amsthm,amssymb}
\usepackage{dsfont,bm,mathtools}
\usepackage[usenames,dvipsnames,svgnames,table]{xcolor}
\usepackage[top=2.3cm,bottom=2.3cm,left=2cm,right=2cm]{geometry}
\usepackage{hyperref}
\allowdisplaybreaks
\emergencystretch=1.5em
\definecolor{ProposalOrange}{RGB}{160,65,0}
\definecolor{RefereePurple}{RGB}{103,34,147}
\newtheorem{lem}{Lemma}[section]
\theoremstyle{remark}
\newtheorem{rem}[lem]{Remark}
\definecolor{ResolutionGreen}{RGB}{0,100,55}
\newcommand{\C}{\mathds{C}}
\newcommand{\Q}{\mathds{Q}}
\newcommand{\rv}{\operatorname{rv}}
\newcommand{\Gal}{\operatorname{Gal}}
\newcommand{\Aut}{\operatorname{Aut}}
\newcommand{\RV}{\mathrm{RV}}
\newcommand{\RES}{\mathrm{RES}}
\newcommand{\VF}{\mathrm{VF}}
\newcommand{\K}{\mathds{k}}
\newcommand{\Z}{\mathds{Z}}
\newcommand{\N}{\mathds{N}}
\newcommand{\vrv}{\operatorname{vrv}}
\newcommand{\GL}{\operatorname{GL}}
\newcommand{\dcl}{\operatorname{dcl}}
\newcommand{\pr}{\operatorname{pr}}
\newcommand{\defn}{\mathrm{def}}
\newcommand{\given}{\,\middle|\,}
\newcommand{\Sbase}{\mathbb S}
\newcommand{\RVlang}{\mathcal L_{\mathrm{RV}}}
\numberwithin{equation}{section}
\pagestyle{plain}
\begin{document}
\setcounter{section}{4}
{\small\noindent Forey--Yin, \emph{Bounded integral and motivic Milnor fiber}.
Section 5 review: second batch (R081--R082).
Approved context is black; proposed manuscript text is orange.}
\medskip
\section{Interlude: RV decomposition and Fubini theorems}\label{csndag}
\setcounter{lem}{11}
\begin{lem}[Monomial decomposition]\label{coh:decom}
Let $U\subseteq(\RV^\times)^n$ be definable. There is a finite definable
partition $(D_i)_i$ of $\vrv(U)$ such that, writing
$U_i=U\cap\vrv^{-1}(D_i)$, there are $M_i\in\GL_n(\Z)$,
$0\leq m_i\leq n$, a nonempty definable set
$V_i\subseteq(\RV^\times)^{m_i}$, and a nonempty definable set
$I_i\subseteq\Gamma^{n-m_i}$ with
\[
 \vrv(V_i)=\{\beta_i\},\qquad
 h_{M_i}(U_i)=V_i\times I_i^\sharp .
\]
All sets and maps are definable over the original base $\Sbase$.
In particular, the partition of $U$ is induced by a partition of
its value support.
\end{lem}
{\normalfont\small\color{ResolutionGreen}
\noindent The proof of Lemma~5.12 has been incorporated as approved.
The next passage replaces the opening of Section~5.1 and its following remark.\par}
\medskip

{\color{ProposalOrange}
\subsection{Tropical motivic Fubini theorem}\label{jp:fubini}
In this subsection we work in the field $\widetilde{\C}$ of complex
Puiseux series, considered as an elementary substructure of
$\mathbb U$. Put $F=\C((t))$ and take the parameter base to be
$\Sbase=F$. Normalize the valuation by $v(t)=1$, so that the
value group of $\widetilde{\C}$ is identified with $\Q$.
All definability and all categories are interpreted in the original
language $\RVlang$ over $\Sbase$. In particular, each
$a/N\in\Q$ is definable over $\Sbase$ as the unique solution of
$N\gamma=a\,v(t)$.

\begin{rem}[The finite Kummer actions]\label{puiseux:kummer:action}
There are natural identifications
\[
 \Gal(\widetilde{\C}/F)\ \cong\ \widehat{\mu}
 \ \cong\ \Aut\bigl(\RV(\widetilde{\C})/\RV(F)\bigr),
 \qquad
 \widehat{\mu}=\varprojlim_N\mu_N,
\]
where the transition map $\mu_M\to\mu_N$ for $N\mid M$
is $\zeta\mapsto\zeta^{M/N}$, and the last automorphism group
uses the induced original-language structure on $\RV$.

For the proof, fix $N\geq1$, choose any root $\tau_N$ of
$X^N-t$ in $\widetilde{\C}$, and put
$F_N=\C((\tau_N))$. Grouping exponents modulo $N$ writes each
element of $F_N$ uniquely as
\[
 \sum_{j=0}^{N-1}\tau_N^j f_j(t),\qquad f_j(t)\in F.
\]
Uniqueness also follows from the pairwise distinct valuations
modulo $\Z$ of the nonzero summands. Thus $[F_N:F]=N$.
Its $F$-automorphisms are precisely
$\tau_N\mapsto\zeta\tau_N$ for $\zeta\in\mu_N$.
Replacing $\tau_N$ by another root multiplies it by an element
of $\mu_N\subseteq F$, so both $F_N$ and the ratio
$\sigma(\tau_N)/\tau_N$ are independent of that choice.
For $N\mid M$ these ratios satisfy
\[
 \frac{\sigma(\tau_N)}{\tau_N}
 =
 \left(\frac{\sigma(\tau_M)}{\tau_M}\right)^{M/N},
\]
since $\tau_M^{M/N}/\tau_N\in\mu_N\subseteq F$.
The fields $F_N$ have union $\widetilde{\C}$.
Their finite automorphisms therefore give exactly the compatible
systems $(\zeta_N)_N\in\widehat{\mu}$, proving the first
identification. Each such automorphism preserves the valuation:
in the displayed expansion it multiplies the $j$th summand by
$\zeta_N^j$, preserving its valuation and the unique least
valuation of a nonzero sum.

Write $r_N=\rv(\tau_N)$ and identify $\C^\times$ with the
value-zero subgroup of $\RV(\widetilde{\C})$.
For any $u$ of value $a/N$ there is a unique $c\in\C^\times$
such that $u=c\,r_N^a$. The induced action is
\[
 \sigma_{(\zeta_N)}(c\,r_N^a)
   =c\,\zeta_N^a r_N^a .
\]
In particular, a nonidentity system moves $r_N$ for some $N$.
The action on $\RV$ is faithful.

Finally, $\RV(F)$ is the subgroup generated by
$\C^\times$ and $\rv(t)$, by taking the leading term of a
Laurent series. An automorphism of the induced RV structure
fixing $\RV(F)$ must send $r_N$ to
$\zeta_Nr_N$ for some $\zeta_N\in\mu_N$: these are precisely
the $N$ roots of $\rv(t)$ in $\RV$.
The same relations as above force compatibility of the
$\zeta_N$. The displayed action then determines the
automorphism on every RV element. It agrees with the unique
valued-field automorphism constructed from that system.
This proves the second identification.
The parameter base for all these definability assertions is
still $\Sbase=F$.
\end{rem}
}
\medskip

{\normalfont\small\color{RefereePurple}
\noindent\textbf{Referee R081 [Clarification; revised] --
Normalization and language in the Puiseux setting.}
For the original computation the referee requests normalization
of the reduced section, compatible roots, and a distinction between
the original and expanded languages so that the chosen roots
do not erase monodromy.\par\medskip
\noindent\textbf{Referee R082 [Typo] -- Formula for a reduced section.}
The original assignment has the wrong input and takes values in
VF instead of RV; the referee requests the corresponding
compatible leading-term formula.\par}
\medskip

{\normalfont\small\color{ProposalOrange}
\noindent\textbf{Proposed disposition of R081--R082.}
Delete the expanded-language Puiseux convention and the
reduced-section parametrization. Retain the Galois interpretation
by the finite Kummer proof above, with the RV action written
explicitly. The roots used in the proof remain auxiliary elements
and are not adjoined to the definability base. No result on
ramification groups or from [19] is used. The construction of
$\Theta$ in the following notation will be reviewed separately
using these finite actions.}
\end{document}
